% 第 19 题：两线圈绕于同一铁芯 + 远处直导线 + 小磁针（按原图重绘）
\begin{tikzpicture}[font=\small, >=Stealth, scale=1.0]
  % 铁芯
  \draw[thick] (0,0) rectangle (5.6,0.5);
  % 左侧线圈（示意为绕圈）
  \foreach \x in {0.35,0.55,...,1.75} { \draw[thick] (\x,0.5) .. controls (\x+0.1,0.95) and (\x+0.1,1.05) .. (\x+0.2,0.5); }
  % 右侧线圈
  \foreach \x in {3.7,3.9,...,5.1} { \draw[thick] (\x,0.5) .. controls (\x+0.1,0.95) and (\x+0.1,1.05) .. (\x+0.2,0.5); }
  % 电源与开关（接右侧线圈）
  \draw[thick] (3.85,0) -- (3.85,-0.85) -- (4.35,-0.85);
  \draw[thick] (4.35,-1.0) -- (4.35,-0.7);
  \draw[thick] (4.55,-1.0) -- (4.55,-0.7);
  \draw[thick] (4.55,-0.85) -- (5.05,-0.85);
  \draw (5.05,-0.85) circle (0.055);
  \draw (5.55,-0.85) circle (0.055);
  \draw[thick] (5.05,-0.79) -- (5.45,-0.55);
  \draw[thick] (5.55,-0.85) -- (5.55,0);
  % 远处直导线（水平）
  \draw[thick] (1.15,2.9) -- (4.35,2.9);
  % 引出导线
  \draw[thick] (0.9,0.5) -- (0.9,1.55) .. controls (0.9,2.1) and (1.3,2.35) .. (1.75,2.4)
      -- (2.05,2.62);
  \draw[thick] (3.0,0.5) -- (3.0,1.55) .. controls (3.0,2.1) and (2.7,2.35) .. (2.3,2.42);
  \draw[thick] (1.15,2.9) .. controls (1.6,2.9) and (1.95,2.8) .. (2.05,2.62);
  \draw[thick] (4.35,2.9) .. controls (3.9,2.9) and (2.55,2.8) .. (2.3,2.42);
  % 小磁针
  \draw[thick] (2.75,2.9) -- (2.75,3.75);
  \draw[thick] (2.3,3.75) -- (3.2,3.75);
  \fill (2.75,3.75) circle (0.06);
  \fill[black] (3.2,3.75) -- (2.95,3.9) -- (2.95,3.6) -- cycle;
  \node[above] at (2.6,3.95) {N};
  \node[above] at (3.15,3.95) {S};
  % 方位标注
  \node[left] at (0.95,2.55) {北};
  \node[right] at (4.5,2.55) {南};
\end{tikzpicture}
